% =========================================================
% ppt/src/minggu-04.tex
% Slide pertemuan 4
% Bab VI — Segmentasi Memori dan Alamat Fisik
%
% Sub-CPMK 1.4 | CPMK01 | CPL02
% =========================================================

\documentclass[aspectratio=169,11pt]{beamer}
% =========================================================
% ppt/src/preamble-ppt.tex
% Tema, warna, paket, dan makro bersama seluruh deck slide.
%
% Dipakai oleh setiap berkas minggu-NN.tex dengan pola:
%   \documentclass[aspectratio=169,11pt]{beamer}
%   % =========================================================
% ppt/src/preamble-ppt.tex
% Tema, warna, paket, dan makro bersama seluruh deck slide.
%
% Dipakai oleh setiap berkas minggu-NN.tex dengan pola:
%   \documentclass[aspectratio=169,11pt]{beamer}
%   % =========================================================
% ppt/src/preamble-ppt.tex
% Tema, warna, paket, dan makro bersama seluruh deck slide.
%
% Dipakai oleh setiap berkas minggu-NN.tex dengan pola:
%   \documentclass[aspectratio=169,11pt]{beamer}
%   \input{preamble-ppt}
%
% Palet warna disalin dari buku_ajar/config/styles.tex agar
% identitas visual slide sama dengan buku ajar.
% =========================================================

% ---------------------------------------------------------
% 1. Paket
% ---------------------------------------------------------
\usepackage[T1]{fontenc}
\usepackage[utf8]{inputenc}
\usepackage[indonesian]{babel}
\usepackage{lmodern}
\usepackage{tikz}
\usetikzlibrary{arrows.meta, positioning, calc, shapes.geometric}
\usepackage{booktabs}
\usepackage{array}
\usepackage{amsmath}
\usepackage{amsfonts}
\usepackage{pifont}          % \ding untuk tanda checklist
\usepackage{listings}

% ---------------------------------------------------------
% 2. Palet warna OBE (identik dengan buku)
% ---------------------------------------------------------
\definecolor{obeBlue}{HTML}{1F3864}
\definecolor{obeBlueLight}{HTML}{DCE4F5}
\definecolor{obeTeal}{HTML}{0F6E6E}
\definecolor{obeTealLight}{HTML}{D9F0EE}
\definecolor{obeGreen}{HTML}{2E6B34}
\definecolor{obeGreenLight}{HTML}{DEEFDF}
\definecolor{obeAmber}{HTML}{8A5A00}
\definecolor{obeAmberLight}{HTML}{FDF1D6}
\definecolor{obeMaroon}{HTML}{7A2233}
\definecolor{obeMaroonLight}{HTML}{F7E1E5}
\definecolor{obeGray}{HTML}{4A4A4A}
\definecolor{obeGrayLight}{HTML}{F0F0F0}
\definecolor{codeBack}{HTML}{F7F7F4}
\definecolor{codeFrame}{HTML}{C9C9C0}

% ---------------------------------------------------------
% 3. Identitas mata kuliah (dipakai seluruh deck)
% ---------------------------------------------------------
\newcommand{\NamaMK}{Bahasa Pemrograman Tingkat Rendah}
\newcommand{\NamaMKShort}{Bhs.\ Pemrograman Tingkat Rendah}
\newcommand{\KodeMK}{MWCA010}
\newcommand{\Prodi}{Program Studi Teknik Informatika}
\newcommand{\Fakultas}{Fakultas Teknologi Informasi}
\newcommand{\PerguruanTinggi}{Universitas Bale Bandung}
\newcommand{\DosenPengampu}{Cecep Suwanda, S.Si., M.Kom.}
\newcommand{\SemesterTeks}{Semester Ganjil}

% \PertemuanTeks dan \JudulBab ditetapkan di masing-masing
% berkas minggu-NN.tex sebelum \begin{document}.
\providecommand{\PertemuanTeks}{Pertemuan}

% ---------------------------------------------------------
% 4. Tema beamer
% ---------------------------------------------------------
\usetheme{default}
\setbeamertemplate{navigation symbols}{}

\setbeamercolor{normal text}{fg=black!88}
\setbeamercolor{structure}{fg=obeBlue}
\setbeamercolor{alerted text}{fg=obeMaroon}

% --- judul slide: pita biru penuh lebar ---
\setbeamercolor{frametitle}{fg=white,bg=obeBlue}
\setbeamerfont{frametitle}{series=\bfseries,size=\large}
\setbeamertemplate{frametitle}{%
    \nointerlineskip%
    \begin{beamercolorbox}[%
        wd=\paperwidth,ht=2.9ex,dp=1.5ex,
        leftskip=0.4cm,rightskip=0.4cm]{frametitle}%
        \usebeamerfont{frametitle}\insertframetitle%
    \end{beamercolorbox}%
    \nointerlineskip%
    \hbox{\textcolor{obeTeal}{\rule{\paperwidth}{1.1pt}}}%
}

% --- kaki slide: identitas mata kuliah, pertemuan, nomor ---
\setbeamercolor{footline}{fg=white,bg=obeGray}
\setbeamerfont{footline}{size=\scriptsize,family=\sffamily}
\setbeamertemplate{footline}{%
    \leavevmode%
    \hbox{%
        \begin{beamercolorbox}[wd=0.42\paperwidth,ht=2.1ex,dp=1.1ex,
            leftskip=0.4cm]{footline}%
            \usebeamerfont{footline}%
            \KodeMK\ \textbullet\ \NamaMKShort%
        \end{beamercolorbox}%
        \begin{beamercolorbox}[wd=0.32\paperwidth,ht=2.1ex,dp=1.1ex,
            center]{footline}%
            \usebeamerfont{footline}\PertemuanTeks%
        \end{beamercolorbox}%
        \begin{beamercolorbox}[wd=0.26\paperwidth,ht=2.1ex,dp=1.1ex,
            rightskip=0.4cm]{footline}%
            \usebeamerfont{footline}%
            \insertframenumber\,/\,\inserttotalframenumber%
        \end{beamercolorbox}%
    }%
    \vskip0pt%
}

% --- daftar ---
\setbeamercolor{itemize item}{fg=obeTeal}
\setbeamercolor{itemize subitem}{fg=obeTeal}
\setbeamercolor{enumerate item}{fg=obeBlue}
\setbeamertemplate{itemize item}{\raise0.2ex\hbox{\tiny$\blacksquare$}}
\setbeamertemplate{itemize subitem}{\raise0.2ex\hbox{\scriptsize$\bullet$}}

% --- kotak (pengganti tcolorbox pada buku) ---
\setbeamercolor{block title}{bg=obeBlueLight,fg=obeBlue}
\setbeamercolor{block body}{bg=obeBlueLight!45,fg=black!88}
\setbeamercolor{block title alerted}{bg=obeMaroonLight,fg=obeMaroon}
\setbeamercolor{block body alerted}{bg=obeMaroonLight!45,fg=black!88}
\setbeamercolor{block title example}{bg=obeTealLight,fg=obeTeal}
\setbeamercolor{block body example}{bg=obeTealLight!45,fg=black!88}

% ---------------------------------------------------------
% 5. Makro notasi (sama persis dengan buku)
%
% Catatan: \note TIDAK didefinisikan di sini karena beamer
% sudah memakai \note untuk catatan pembicara. Untuk catatan
% di slide dipakai \catatan.
% ---------------------------------------------------------
\newcommand{\code}[1]{\texttt{#1}}
\newcommand{\asm}[1]{\texttt{#1}}
\newcommand{\reg}[1]{\texttt{#1}}
\newcommand{\flag}[1]{\texttt{#1}}
\newcommand{\dir}[1]{\texttt{#1}}
\newcommand{\file}[1]{\texttt{#1}}
\newcommand{\eng}[1]{\textit{#1}}

\newcommand{\hex}[1]{\texttt{#1h}}
\newcommand{\bin}[1]{\texttt{#1b}}
\newcommand{\dec}[1]{\texttt{#1}}

\newcommand{\important}[1]{\textbf{Penting:} #1}
\newcommand{\catatan}[1]{\textbf{Catatan:} #1}
\newcommand{\term}[1]{\textbf{#1}}

\newcommand{\CPL}[1]{\textbf{CPL#1}}
\newcommand{\CPMK}[1]{\textbf{CPMK#1}}
\newcommand{\SubCPMK}[1]{\textbf{Sub-CPMK #1}}

\newcommand{\checkbox}{\ding{113}\hspace{0.35em}}
\newcommand{\checked}{\ding{51}\hspace{0.35em}}
\newcommand{\crossed}{\ding{55}\hspace{0.35em}}

\newcommand{\flowarrow}{\ensuremath{\rightarrow}}

% ---------------------------------------------------------
% 6. Makro tata letak slide
% ---------------------------------------------------------

% Keterangan sumber materi (ketertelusuran ke buku)
\newcommand{\Sumber}[1]{%
    \vfill
    {\scriptsize\color{obeGray}Sumber: #1\par}%
}

% Frame judul deck: \JudulDeck{<judul bab>}{<pertemuan>}
\newcommand{\JudulDeck}[2]{%
\begin{frame}[plain]
    \vspace{0.25cm}
    {\color{obeBlue}\sffamily\bfseries\small \NamaMK\ (\KodeMK)\par}
    \vspace{0.12cm}
    {\color{obeGray}\sffamily\scriptsize
     \Prodi\ \textbullet\ \Fakultas\ \textbullet\ \PerguruanTinggi\par}
    \vfill
    {\color{obeBlue}\sffamily\LARGE\bfseries #1\par}
    \vspace{0.4cm}
    {\color{obeTeal}\sffamily\large #2\par}
    \vspace{0.45cm}
    {\color{obeTeal}\rule{\textwidth}{1pt}}
    \vfill
    {\color{obeGray}\sffamily\small Dosen Pengampu\par}
    \vspace{0.08cm}
    {\sffamily\bfseries \DosenPengampu\par}
    \vspace{0.2cm}
    {\color{obeGray}\sffamily\small \SemesterTeks\par}
    \vspace{0.25cm}
\end{frame}
}

% Frame pemisah bagian
\newcommand{\PemisahBagian}[1]{%
\begin{frame}[plain]
    \vfill
    \begin{center}
        {\color{obeTeal}\sffamily\small\bfseries BAGIAN\par}
        \vspace{0.35cm}
        {\color{obeBlue}\sffamily\Large\bfseries #1\par}
        \vspace{0.35cm}
        {\color{obeTeal}\rule{0.55\textwidth}{0.9pt}}
    \end{center}
    \vfill
\end{frame}
}

% ---------------------------------------------------------
% 7. Gaya penampil kode
% ---------------------------------------------------------
\input{listing-ppt}

%
% Palet warna disalin dari buku_ajar/config/styles.tex agar
% identitas visual slide sama dengan buku ajar.
% =========================================================

% ---------------------------------------------------------
% 1. Paket
% ---------------------------------------------------------
\usepackage[T1]{fontenc}
\usepackage[utf8]{inputenc}
\usepackage[indonesian]{babel}
\usepackage{lmodern}
\usepackage{tikz}
\usetikzlibrary{arrows.meta, positioning, calc, shapes.geometric}
\usepackage{booktabs}
\usepackage{array}
\usepackage{amsmath}
\usepackage{amsfonts}
\usepackage{pifont}          % \ding untuk tanda checklist
\usepackage{listings}

% ---------------------------------------------------------
% 2. Palet warna OBE (identik dengan buku)
% ---------------------------------------------------------
\definecolor{obeBlue}{HTML}{1F3864}
\definecolor{obeBlueLight}{HTML}{DCE4F5}
\definecolor{obeTeal}{HTML}{0F6E6E}
\definecolor{obeTealLight}{HTML}{D9F0EE}
\definecolor{obeGreen}{HTML}{2E6B34}
\definecolor{obeGreenLight}{HTML}{DEEFDF}
\definecolor{obeAmber}{HTML}{8A5A00}
\definecolor{obeAmberLight}{HTML}{FDF1D6}
\definecolor{obeMaroon}{HTML}{7A2233}
\definecolor{obeMaroonLight}{HTML}{F7E1E5}
\definecolor{obeGray}{HTML}{4A4A4A}
\definecolor{obeGrayLight}{HTML}{F0F0F0}
\definecolor{codeBack}{HTML}{F7F7F4}
\definecolor{codeFrame}{HTML}{C9C9C0}

% ---------------------------------------------------------
% 3. Identitas mata kuliah (dipakai seluruh deck)
% ---------------------------------------------------------
\newcommand{\NamaMK}{Bahasa Pemrograman Tingkat Rendah}
\newcommand{\NamaMKShort}{Bhs.\ Pemrograman Tingkat Rendah}
\newcommand{\KodeMK}{MWCA010}
\newcommand{\Prodi}{Program Studi Teknik Informatika}
\newcommand{\Fakultas}{Fakultas Teknologi Informasi}
\newcommand{\PerguruanTinggi}{Universitas Bale Bandung}
\newcommand{\DosenPengampu}{Cecep Suwanda, S.Si., M.Kom.}
\newcommand{\SemesterTeks}{Semester Ganjil}

% \PertemuanTeks dan \JudulBab ditetapkan di masing-masing
% berkas minggu-NN.tex sebelum \begin{document}.
\providecommand{\PertemuanTeks}{Pertemuan}

% ---------------------------------------------------------
% 4. Tema beamer
% ---------------------------------------------------------
\usetheme{default}
\setbeamertemplate{navigation symbols}{}

\setbeamercolor{normal text}{fg=black!88}
\setbeamercolor{structure}{fg=obeBlue}
\setbeamercolor{alerted text}{fg=obeMaroon}

% --- judul slide: pita biru penuh lebar ---
\setbeamercolor{frametitle}{fg=white,bg=obeBlue}
\setbeamerfont{frametitle}{series=\bfseries,size=\large}
\setbeamertemplate{frametitle}{%
    \nointerlineskip%
    \begin{beamercolorbox}[%
        wd=\paperwidth,ht=2.9ex,dp=1.5ex,
        leftskip=0.4cm,rightskip=0.4cm]{frametitle}%
        \usebeamerfont{frametitle}\insertframetitle%
    \end{beamercolorbox}%
    \nointerlineskip%
    \hbox{\textcolor{obeTeal}{\rule{\paperwidth}{1.1pt}}}%
}

% --- kaki slide: identitas mata kuliah, pertemuan, nomor ---
\setbeamercolor{footline}{fg=white,bg=obeGray}
\setbeamerfont{footline}{size=\scriptsize,family=\sffamily}
\setbeamertemplate{footline}{%
    \leavevmode%
    \hbox{%
        \begin{beamercolorbox}[wd=0.42\paperwidth,ht=2.1ex,dp=1.1ex,
            leftskip=0.4cm]{footline}%
            \usebeamerfont{footline}%
            \KodeMK\ \textbullet\ \NamaMKShort%
        \end{beamercolorbox}%
        \begin{beamercolorbox}[wd=0.32\paperwidth,ht=2.1ex,dp=1.1ex,
            center]{footline}%
            \usebeamerfont{footline}\PertemuanTeks%
        \end{beamercolorbox}%
        \begin{beamercolorbox}[wd=0.26\paperwidth,ht=2.1ex,dp=1.1ex,
            rightskip=0.4cm]{footline}%
            \usebeamerfont{footline}%
            \insertframenumber\,/\,\inserttotalframenumber%
        \end{beamercolorbox}%
    }%
    \vskip0pt%
}

% --- daftar ---
\setbeamercolor{itemize item}{fg=obeTeal}
\setbeamercolor{itemize subitem}{fg=obeTeal}
\setbeamercolor{enumerate item}{fg=obeBlue}
\setbeamertemplate{itemize item}{\raise0.2ex\hbox{\tiny$\blacksquare$}}
\setbeamertemplate{itemize subitem}{\raise0.2ex\hbox{\scriptsize$\bullet$}}

% --- kotak (pengganti tcolorbox pada buku) ---
\setbeamercolor{block title}{bg=obeBlueLight,fg=obeBlue}
\setbeamercolor{block body}{bg=obeBlueLight!45,fg=black!88}
\setbeamercolor{block title alerted}{bg=obeMaroonLight,fg=obeMaroon}
\setbeamercolor{block body alerted}{bg=obeMaroonLight!45,fg=black!88}
\setbeamercolor{block title example}{bg=obeTealLight,fg=obeTeal}
\setbeamercolor{block body example}{bg=obeTealLight!45,fg=black!88}

% ---------------------------------------------------------
% 5. Makro notasi (sama persis dengan buku)
%
% Catatan: \note TIDAK didefinisikan di sini karena beamer
% sudah memakai \note untuk catatan pembicara. Untuk catatan
% di slide dipakai \catatan.
% ---------------------------------------------------------
\newcommand{\code}[1]{\texttt{#1}}
\newcommand{\asm}[1]{\texttt{#1}}
\newcommand{\reg}[1]{\texttt{#1}}
\newcommand{\flag}[1]{\texttt{#1}}
\newcommand{\dir}[1]{\texttt{#1}}
\newcommand{\file}[1]{\texttt{#1}}
\newcommand{\eng}[1]{\textit{#1}}

\newcommand{\hex}[1]{\texttt{#1h}}
\newcommand{\bin}[1]{\texttt{#1b}}
\newcommand{\dec}[1]{\texttt{#1}}

\newcommand{\important}[1]{\textbf{Penting:} #1}
\newcommand{\catatan}[1]{\textbf{Catatan:} #1}
\newcommand{\term}[1]{\textbf{#1}}

\newcommand{\CPL}[1]{\textbf{CPL#1}}
\newcommand{\CPMK}[1]{\textbf{CPMK#1}}
\newcommand{\SubCPMK}[1]{\textbf{Sub-CPMK #1}}

\newcommand{\checkbox}{\ding{113}\hspace{0.35em}}
\newcommand{\checked}{\ding{51}\hspace{0.35em}}
\newcommand{\crossed}{\ding{55}\hspace{0.35em}}

\newcommand{\flowarrow}{\ensuremath{\rightarrow}}

% ---------------------------------------------------------
% 6. Makro tata letak slide
% ---------------------------------------------------------

% Keterangan sumber materi (ketertelusuran ke buku)
\newcommand{\Sumber}[1]{%
    \vfill
    {\scriptsize\color{obeGray}Sumber: #1\par}%
}

% Frame judul deck: \JudulDeck{<judul bab>}{<pertemuan>}
\newcommand{\JudulDeck}[2]{%
\begin{frame}[plain]
    \vspace{0.25cm}
    {\color{obeBlue}\sffamily\bfseries\small \NamaMK\ (\KodeMK)\par}
    \vspace{0.12cm}
    {\color{obeGray}\sffamily\scriptsize
     \Prodi\ \textbullet\ \Fakultas\ \textbullet\ \PerguruanTinggi\par}
    \vfill
    {\color{obeBlue}\sffamily\LARGE\bfseries #1\par}
    \vspace{0.4cm}
    {\color{obeTeal}\sffamily\large #2\par}
    \vspace{0.45cm}
    {\color{obeTeal}\rule{\textwidth}{1pt}}
    \vfill
    {\color{obeGray}\sffamily\small Dosen Pengampu\par}
    \vspace{0.08cm}
    {\sffamily\bfseries \DosenPengampu\par}
    \vspace{0.2cm}
    {\color{obeGray}\sffamily\small \SemesterTeks\par}
    \vspace{0.25cm}
\end{frame}
}

% Frame pemisah bagian
\newcommand{\PemisahBagian}[1]{%
\begin{frame}[plain]
    \vfill
    \begin{center}
        {\color{obeTeal}\sffamily\small\bfseries BAGIAN\par}
        \vspace{0.35cm}
        {\color{obeBlue}\sffamily\Large\bfseries #1\par}
        \vspace{0.35cm}
        {\color{obeTeal}\rule{0.55\textwidth}{0.9pt}}
    \end{center}
    \vfill
\end{frame}
}

% ---------------------------------------------------------
% 7. Gaya penampil kode
% ---------------------------------------------------------
% =========================================================
% ppt/src/listing-ppt.tex
% Gaya penampil kode untuk slide.
%
% Diturunkan dari buku_ajar/config/listings.tex, tetapi
% disesuaikan untuk slide: ukuran lebih kecil (karena frame
% 16:9 dan kode harus utuh terlihat) dan jarak lebih rapat.
% =========================================================

% ---------------------------------------------------------
% 1. Definisi bahasa MASM/TASM 8086 (salinan dari buku)
% ---------------------------------------------------------
\lstdefinelanguage{MASM}{
    morekeywords={
        MOV, XCHG, LEA, LDS, LES, PUSH, POP, PUSHA, POPA, PUSHF, POPF,
        LAHF, SAHF, XLAT, IN, OUT, CBW, CWD,
        ADD, ADC, SUB, SBB, INC, DEC, NEG, CMP, MUL, IMUL, DIV, IDIV,
        AAA, AAS, AAM, AAD, DAA, DAS,
        AND, OR, XOR, NOT, TEST, SHL, SHR, SAL, SAR, ROL, ROR, RCL, RCR,
        JMP, JE, JZ, JNE, JNZ, JG, JGE, JL, JLE, JA, JAE, JB, JBE,
        JC, JNC, JS, JNS, JO, JNO, JP, JPE, JNP, JPO, JCXZ,
        LOOP, LOOPE, LOOPZ, LOOPNE, LOOPNZ, CALL, RET, INT, IRET, INTO,
        HLT, NOP, CLI, STI, CLD, STD, CLC, STC, CMC, WAIT, ESC, LOCK,
        DB, DW, DD, DQ, DT, EQU, ORG, SEGMENT, ENDS, ASSUME, END,
        PROC, ENDP, MACRO, ENDM, LOCAL, PUBLIC, EXTRN, EXTERN,
        INCLUDE, INCLUDELIB, MODEL, STACK, CODE, DATA, CONST, ALIGN, EVEN,
        OFFSET, PTR, BYTE, WORD, DWORD, QWORD, TBYTE,
        NEAR, FAR, SHORT, LABEL, TYPE, SIZE, LENGTH, DUP, THIS,
        AX, BX, CX, DX, AL, AH, BL, BH, CL, CH, DL, DH,
        SP, BP, SI, DI, IP, CS, DS, SS, ES, FLAGS,
        CF, PF, AF, ZF, TF, IF, DF, OF
    },
    sensitive=false,
    morecomment=[l]{;},
    morestring=[b]',
    morestring=[b]"
}

% ---------------------------------------------------------
% 2. Gaya kode assembly untuk slide
% ---------------------------------------------------------
\lstdefinestyle{asmslide}{
    language=MASM,
    basicstyle=\ttfamily\scriptsize,
    keywordstyle=\color{obeBlue}\bfseries,
    commentstyle=\color{obeGreen}\itshape,
    stringstyle=\color{obeMaroon},
    numberstyle=\tiny\color{obeGray},
    numbers=left,
    numbersep=5pt,
    stepnumber=1,
    backgroundcolor=\color{codeBack},
    frame=single,
    rulecolor=\color{codeFrame},
    framesep=4pt,
    breaklines=true,
    breakatwhitespace=false,
    % CATATAN PENTING: opsi `postbreak` TIDAK boleh dipakai di sini.
    % Di beamer, `breaklines=true` bersama `postbreak=<kode>` membuat
    % listings gagal dirakit dengan galat
    %   "Missing number, treated as zero <to be read again> \kern".
    % Karena \lstset{style=asmslide} bersifat global, opsi postbreak
    % juga bocor ke gaya lain (mis. outslide/keluaran) dan merusaknya.
    % Di buku ajar (kelas `book`) justru sebaliknya: postbreak aman dan
    % dipakai. Jadi larangan ini khusus untuk slide.
    showstringspaces=false,
    tabsize=4,
    columns=fullflexible,
    keepspaces=true,
    upquote=true,
    xleftmargin=12pt,
    aboveskip=0.4em,
    belowskip=0.4em
}

% ---------------------------------------------------------
% 3. Gaya keluaran program / sesi debugger untuk slide
% ---------------------------------------------------------
\lstdefinestyle{outslide}{
    basicstyle=\ttfamily\scriptsize,
    keywordstyle=\color{black},
    commentstyle=\color{obeGray},
    stringstyle=\color{black},
    numbers=none,
    backgroundcolor=\color{obeGrayLight},
    frame=single,
    rulecolor=\color{codeFrame},
    framesep=4pt,
    breaklines=true,
    showstringspaces=false,
    columns=fullflexible,
    keepspaces=true,
    upquote=true,
    xleftmargin=4pt,
    aboveskip=0.4em,
    belowskip=0.4em
}

\lstset{style=asmslide}

% ---------------------------------------------------------
% 4. Environment ringkas untuk dipakai di dalam frame
%
% Frame yang memuat kedua environment ini WAJIB memakai
% opsi [fragile], karena isinya verbatim.
% ---------------------------------------------------------
\lstnewenvironment{kodeasm}[1][]{%
    \lstset{style=asmslide, #1}%
}{}

\lstnewenvironment{keluaran}[1][]{%
    \lstset{style=outslide, #1}%
}{}


%
% Palet warna disalin dari buku_ajar/config/styles.tex agar
% identitas visual slide sama dengan buku ajar.
% =========================================================

% ---------------------------------------------------------
% 1. Paket
% ---------------------------------------------------------
\usepackage[T1]{fontenc}
\usepackage[utf8]{inputenc}
\usepackage[indonesian]{babel}
\usepackage{lmodern}
\usepackage{tikz}
\usetikzlibrary{arrows.meta, positioning, calc, shapes.geometric}
\usepackage{booktabs}
\usepackage{array}
\usepackage{amsmath}
\usepackage{amsfonts}
\usepackage{pifont}          % \ding untuk tanda checklist
\usepackage{listings}

% ---------------------------------------------------------
% 2. Palet warna OBE (identik dengan buku)
% ---------------------------------------------------------
\definecolor{obeBlue}{HTML}{1F3864}
\definecolor{obeBlueLight}{HTML}{DCE4F5}
\definecolor{obeTeal}{HTML}{0F6E6E}
\definecolor{obeTealLight}{HTML}{D9F0EE}
\definecolor{obeGreen}{HTML}{2E6B34}
\definecolor{obeGreenLight}{HTML}{DEEFDF}
\definecolor{obeAmber}{HTML}{8A5A00}
\definecolor{obeAmberLight}{HTML}{FDF1D6}
\definecolor{obeMaroon}{HTML}{7A2233}
\definecolor{obeMaroonLight}{HTML}{F7E1E5}
\definecolor{obeGray}{HTML}{4A4A4A}
\definecolor{obeGrayLight}{HTML}{F0F0F0}
\definecolor{codeBack}{HTML}{F7F7F4}
\definecolor{codeFrame}{HTML}{C9C9C0}

% ---------------------------------------------------------
% 3. Identitas mata kuliah (dipakai seluruh deck)
% ---------------------------------------------------------
\newcommand{\NamaMK}{Bahasa Pemrograman Tingkat Rendah}
\newcommand{\NamaMKShort}{Bhs.\ Pemrograman Tingkat Rendah}
\newcommand{\KodeMK}{MWCA010}
\newcommand{\Prodi}{Program Studi Teknik Informatika}
\newcommand{\Fakultas}{Fakultas Teknologi Informasi}
\newcommand{\PerguruanTinggi}{Universitas Bale Bandung}
\newcommand{\DosenPengampu}{Cecep Suwanda, S.Si., M.Kom.}
\newcommand{\SemesterTeks}{Semester Ganjil}

% \PertemuanTeks dan \JudulBab ditetapkan di masing-masing
% berkas minggu-NN.tex sebelum \begin{document}.
\providecommand{\PertemuanTeks}{Pertemuan}

% ---------------------------------------------------------
% 4. Tema beamer
% ---------------------------------------------------------
\usetheme{default}
\setbeamertemplate{navigation symbols}{}

\setbeamercolor{normal text}{fg=black!88}
\setbeamercolor{structure}{fg=obeBlue}
\setbeamercolor{alerted text}{fg=obeMaroon}

% --- judul slide: pita biru penuh lebar ---
\setbeamercolor{frametitle}{fg=white,bg=obeBlue}
\setbeamerfont{frametitle}{series=\bfseries,size=\large}
\setbeamertemplate{frametitle}{%
    \nointerlineskip%
    \begin{beamercolorbox}[%
        wd=\paperwidth,ht=2.9ex,dp=1.5ex,
        leftskip=0.4cm,rightskip=0.4cm]{frametitle}%
        \usebeamerfont{frametitle}\insertframetitle%
    \end{beamercolorbox}%
    \nointerlineskip%
    \hbox{\textcolor{obeTeal}{\rule{\paperwidth}{1.1pt}}}%
}

% --- kaki slide: identitas mata kuliah, pertemuan, nomor ---
\setbeamercolor{footline}{fg=white,bg=obeGray}
\setbeamerfont{footline}{size=\scriptsize,family=\sffamily}
\setbeamertemplate{footline}{%
    \leavevmode%
    \hbox{%
        \begin{beamercolorbox}[wd=0.42\paperwidth,ht=2.1ex,dp=1.1ex,
            leftskip=0.4cm]{footline}%
            \usebeamerfont{footline}%
            \KodeMK\ \textbullet\ \NamaMKShort%
        \end{beamercolorbox}%
        \begin{beamercolorbox}[wd=0.32\paperwidth,ht=2.1ex,dp=1.1ex,
            center]{footline}%
            \usebeamerfont{footline}\PertemuanTeks%
        \end{beamercolorbox}%
        \begin{beamercolorbox}[wd=0.26\paperwidth,ht=2.1ex,dp=1.1ex,
            rightskip=0.4cm]{footline}%
            \usebeamerfont{footline}%
            \insertframenumber\,/\,\inserttotalframenumber%
        \end{beamercolorbox}%
    }%
    \vskip0pt%
}

% --- daftar ---
\setbeamercolor{itemize item}{fg=obeTeal}
\setbeamercolor{itemize subitem}{fg=obeTeal}
\setbeamercolor{enumerate item}{fg=obeBlue}
\setbeamertemplate{itemize item}{\raise0.2ex\hbox{\tiny$\blacksquare$}}
\setbeamertemplate{itemize subitem}{\raise0.2ex\hbox{\scriptsize$\bullet$}}

% --- kotak (pengganti tcolorbox pada buku) ---
\setbeamercolor{block title}{bg=obeBlueLight,fg=obeBlue}
\setbeamercolor{block body}{bg=obeBlueLight!45,fg=black!88}
\setbeamercolor{block title alerted}{bg=obeMaroonLight,fg=obeMaroon}
\setbeamercolor{block body alerted}{bg=obeMaroonLight!45,fg=black!88}
\setbeamercolor{block title example}{bg=obeTealLight,fg=obeTeal}
\setbeamercolor{block body example}{bg=obeTealLight!45,fg=black!88}

% ---------------------------------------------------------
% 5. Makro notasi (sama persis dengan buku)
%
% Catatan: \note TIDAK didefinisikan di sini karena beamer
% sudah memakai \note untuk catatan pembicara. Untuk catatan
% di slide dipakai \catatan.
% ---------------------------------------------------------
\newcommand{\code}[1]{\texttt{#1}}
\newcommand{\asm}[1]{\texttt{#1}}
\newcommand{\reg}[1]{\texttt{#1}}
\newcommand{\flag}[1]{\texttt{#1}}
\newcommand{\dir}[1]{\texttt{#1}}
\newcommand{\file}[1]{\texttt{#1}}
\newcommand{\eng}[1]{\textit{#1}}

\newcommand{\hex}[1]{\texttt{#1h}}
\newcommand{\bin}[1]{\texttt{#1b}}
\newcommand{\dec}[1]{\texttt{#1}}

\newcommand{\important}[1]{\textbf{Penting:} #1}
\newcommand{\catatan}[1]{\textbf{Catatan:} #1}
\newcommand{\term}[1]{\textbf{#1}}

\newcommand{\CPL}[1]{\textbf{CPL#1}}
\newcommand{\CPMK}[1]{\textbf{CPMK#1}}
\newcommand{\SubCPMK}[1]{\textbf{Sub-CPMK #1}}

\newcommand{\checkbox}{\ding{113}\hspace{0.35em}}
\newcommand{\checked}{\ding{51}\hspace{0.35em}}
\newcommand{\crossed}{\ding{55}\hspace{0.35em}}

\newcommand{\flowarrow}{\ensuremath{\rightarrow}}

% ---------------------------------------------------------
% 6. Makro tata letak slide
% ---------------------------------------------------------

% Keterangan sumber materi (ketertelusuran ke buku)
\newcommand{\Sumber}[1]{%
    \vfill
    {\scriptsize\color{obeGray}Sumber: #1\par}%
}

% Frame judul deck: \JudulDeck{<judul bab>}{<pertemuan>}
\newcommand{\JudulDeck}[2]{%
\begin{frame}[plain]
    \vspace{0.25cm}
    {\color{obeBlue}\sffamily\bfseries\small \NamaMK\ (\KodeMK)\par}
    \vspace{0.12cm}
    {\color{obeGray}\sffamily\scriptsize
     \Prodi\ \textbullet\ \Fakultas\ \textbullet\ \PerguruanTinggi\par}
    \vfill
    {\color{obeBlue}\sffamily\LARGE\bfseries #1\par}
    \vspace{0.4cm}
    {\color{obeTeal}\sffamily\large #2\par}
    \vspace{0.45cm}
    {\color{obeTeal}\rule{\textwidth}{1pt}}
    \vfill
    {\color{obeGray}\sffamily\small Dosen Pengampu\par}
    \vspace{0.08cm}
    {\sffamily\bfseries \DosenPengampu\par}
    \vspace{0.2cm}
    {\color{obeGray}\sffamily\small \SemesterTeks\par}
    \vspace{0.25cm}
\end{frame}
}

% Frame pemisah bagian
\newcommand{\PemisahBagian}[1]{%
\begin{frame}[plain]
    \vfill
    \begin{center}
        {\color{obeTeal}\sffamily\small\bfseries BAGIAN\par}
        \vspace{0.35cm}
        {\color{obeBlue}\sffamily\Large\bfseries #1\par}
        \vspace{0.35cm}
        {\color{obeTeal}\rule{0.55\textwidth}{0.9pt}}
    \end{center}
    \vfill
\end{frame}
}

% ---------------------------------------------------------
% 7. Gaya penampil kode
% ---------------------------------------------------------
% =========================================================
% ppt/src/listing-ppt.tex
% Gaya penampil kode untuk slide.
%
% Diturunkan dari buku_ajar/config/listings.tex, tetapi
% disesuaikan untuk slide: ukuran lebih kecil (karena frame
% 16:9 dan kode harus utuh terlihat) dan jarak lebih rapat.
% =========================================================

% ---------------------------------------------------------
% 1. Definisi bahasa MASM/TASM 8086 (salinan dari buku)
% ---------------------------------------------------------
\lstdefinelanguage{MASM}{
    morekeywords={
        MOV, XCHG, LEA, LDS, LES, PUSH, POP, PUSHA, POPA, PUSHF, POPF,
        LAHF, SAHF, XLAT, IN, OUT, CBW, CWD,
        ADD, ADC, SUB, SBB, INC, DEC, NEG, CMP, MUL, IMUL, DIV, IDIV,
        AAA, AAS, AAM, AAD, DAA, DAS,
        AND, OR, XOR, NOT, TEST, SHL, SHR, SAL, SAR, ROL, ROR, RCL, RCR,
        JMP, JE, JZ, JNE, JNZ, JG, JGE, JL, JLE, JA, JAE, JB, JBE,
        JC, JNC, JS, JNS, JO, JNO, JP, JPE, JNP, JPO, JCXZ,
        LOOP, LOOPE, LOOPZ, LOOPNE, LOOPNZ, CALL, RET, INT, IRET, INTO,
        HLT, NOP, CLI, STI, CLD, STD, CLC, STC, CMC, WAIT, ESC, LOCK,
        DB, DW, DD, DQ, DT, EQU, ORG, SEGMENT, ENDS, ASSUME, END,
        PROC, ENDP, MACRO, ENDM, LOCAL, PUBLIC, EXTRN, EXTERN,
        INCLUDE, INCLUDELIB, MODEL, STACK, CODE, DATA, CONST, ALIGN, EVEN,
        OFFSET, PTR, BYTE, WORD, DWORD, QWORD, TBYTE,
        NEAR, FAR, SHORT, LABEL, TYPE, SIZE, LENGTH, DUP, THIS,
        AX, BX, CX, DX, AL, AH, BL, BH, CL, CH, DL, DH,
        SP, BP, SI, DI, IP, CS, DS, SS, ES, FLAGS,
        CF, PF, AF, ZF, TF, IF, DF, OF
    },
    sensitive=false,
    morecomment=[l]{;},
    morestring=[b]',
    morestring=[b]"
}

% ---------------------------------------------------------
% 2. Gaya kode assembly untuk slide
% ---------------------------------------------------------
\lstdefinestyle{asmslide}{
    language=MASM,
    basicstyle=\ttfamily\scriptsize,
    keywordstyle=\color{obeBlue}\bfseries,
    commentstyle=\color{obeGreen}\itshape,
    stringstyle=\color{obeMaroon},
    numberstyle=\tiny\color{obeGray},
    numbers=left,
    numbersep=5pt,
    stepnumber=1,
    backgroundcolor=\color{codeBack},
    frame=single,
    rulecolor=\color{codeFrame},
    framesep=4pt,
    breaklines=true,
    breakatwhitespace=false,
    % CATATAN PENTING: opsi `postbreak` TIDAK boleh dipakai di sini.
    % Di beamer, `breaklines=true` bersama `postbreak=<kode>` membuat
    % listings gagal dirakit dengan galat
    %   "Missing number, treated as zero <to be read again> \kern".
    % Karena \lstset{style=asmslide} bersifat global, opsi postbreak
    % juga bocor ke gaya lain (mis. outslide/keluaran) dan merusaknya.
    % Di buku ajar (kelas `book`) justru sebaliknya: postbreak aman dan
    % dipakai. Jadi larangan ini khusus untuk slide.
    showstringspaces=false,
    tabsize=4,
    columns=fullflexible,
    keepspaces=true,
    upquote=true,
    xleftmargin=12pt,
    aboveskip=0.4em,
    belowskip=0.4em
}

% ---------------------------------------------------------
% 3. Gaya keluaran program / sesi debugger untuk slide
% ---------------------------------------------------------
\lstdefinestyle{outslide}{
    basicstyle=\ttfamily\scriptsize,
    keywordstyle=\color{black},
    commentstyle=\color{obeGray},
    stringstyle=\color{black},
    numbers=none,
    backgroundcolor=\color{obeGrayLight},
    frame=single,
    rulecolor=\color{codeFrame},
    framesep=4pt,
    breaklines=true,
    showstringspaces=false,
    columns=fullflexible,
    keepspaces=true,
    upquote=true,
    xleftmargin=4pt,
    aboveskip=0.4em,
    belowskip=0.4em
}

\lstset{style=asmslide}

% ---------------------------------------------------------
% 4. Environment ringkas untuk dipakai di dalam frame
%
% Frame yang memuat kedua environment ini WAJIB memakai
% opsi [fragile], karena isinya verbatim.
% ---------------------------------------------------------
\lstnewenvironment{kodeasm}[1][]{%
    \lstset{style=asmslide, #1}%
}{}

\lstnewenvironment{keluaran}[1][]{%
    \lstset{style=outslide, #1}%
}{}



\renewcommand{\PertemuanTeks}{Pertemuan 4}

\begin{document}

% =========================================================
% 1. JUDUL
% =========================================================
\JudulDeck{Segmentasi Memori dan\\ Alamat Fisik}{Pertemuan 4}

% =========================================================
% 2. CAPAIAN DAN TUJUAN
% =========================================================
\begin{frame}{Capaian dan Tujuan Pembelajaran}
\begin{block}{Sub-CPMK 1.4}
\small Menghitung alamat fisik dari pasangan
\eng{segment:offset} serta menjelaskan pembagian memori 1 MB dan
\eng{overlapping} segmen.
\end{block}

\vspace{0.1em}
{\small\textbf{CPMK yang didukung:} \CPMK{01}\quad\textbullet\quad
\textbf{CPL yang didukung:} \CPL{02}\par}

\vspace{0.15em}
\begin{exampleblock}{Setelah pertemuan ini, mahasiswa mampu:}
\footnotesize
\begin{enumerate}
    \item menjelaskan alasan 8086 memakai register segmen 16 bit
          walaupun bus alamatnya 20 bit;
    \item membedakan alamat logis dan alamat fisik;
    \item menghitung alamat fisik dari pasangan
          \eng{segment:offset};
    \item menjelaskan paragraf 16 byte, segmen 64 KB, dan
          \eng{overlapping};
    \item menyebutkan pembagian memori 1 MB serta memverifikasi
          alamat dengan DEBUG.
\end{enumerate}
\end{exampleblock}
\end{frame}

% =========================================================
% 3. PETA MATERI
% =========================================================
\begin{frame}{Peta Materi}
\small
\begin{enumerate}
    \item \textbf{Mengapa Segmentasi Diperlukan}\\
          {\footnotesize register 16 bit dan bus 20 bit; alamat
          logis dan fisik; rumus; perhitungan bertingkat}

    \vspace{0.25em}
    \item \textbf{Paragraf 16 Byte, Segmen 64 KB, dan Overlapping}\\
          {\footnotesize paragraf; jendela 64 KB; satu alamat banyak
          pasangan; operator OFFSET dan SEG}

    \vspace{0.25em}
    \item \textbf{Pembagian Memori 1 MB dan Verifikasi Alamat}\\
          {\footnotesize enam belas blok; peta memori; DEBUG;
          pembuktian \eng{overlapping}}
\end{enumerate}

\vspace{0.35em}
\begin{block}{Praktik pada pertemuan ini}
\small Menghitung alamat fisik, memverifikasinya dengan DEBUG, dan
membuktikan \eng{overlapping} segmen.
\end{block}
\end{frame}

% =========================================================
% BAGIAN 1
% =========================================================
\PemisahBagian{Mengapa Segmentasi Diperlukan}

% ---------------------------------------------------------
\begin{frame}{Register 16 Bit vs Bus Alamat 20 Bit}
\footnotesize
\begin{center}
\begin{tabular}{@{}>{\raggedright\arraybackslash}p{4.0cm}
    >{\raggedright\arraybackslash}p{1.8cm}
    >{\raggedright\arraybackslash}p{2.4cm}
    >{\raggedright\arraybackslash}p{4.2cm}@{}}
    \toprule
    \textbf{Sisi} & \textbf{Lebar} & \textbf{Jangkauan} &
    \textbf{Akibat} \\
    \midrule
    Register 16 bit & 16 bit & 65.536 alamat &
    Hanya 64 KB yang dapat ditunjuk langsung \\
    \addlinespace
    Bus alamat prosesor & 20 bit & 1.048.576 alamat &
    Seluruh 1 MB dapat dijangkau \\
    \bottomrule
\end{tabular}
\end{center}

\vspace{0.3em}
{\color{obeGray}Ada selisih jangkauan 16 kali. 8086 menjembatani
keduanya dengan pasangan \eng{segment:offset}.\par}
\Sumber{Bab VI, Bagian 1}
\end{frame}

% ---------------------------------------------------------
\begin{frame}{Alamat Logis dan Alamat Fisik}
\small
\begin{block}{Alamat logis}
\small Cara programmer menuliskan alamat: pasangan
\eng{segment:offset}, misalnya \code{1234h:5678h}. Nilai
\hex{1234} adalah segmen dan \hex{5678} adalah offset.
\end{block}

\vspace{0.2em}
\begin{block}{Alamat fisik}
\small Alamat 20 bit yang dibangkitkan prosesor dan dikirim ke bus
alamat untuk memilih satu lokasi memori.
\end{block}

\vspace{0.2em}
\begin{alertblock}{Nilai segmen bukan alamat}
\small Segmen \hex{1234} tidak menunjuk alamat \hex{1234},
melainkan awal paragraf pada \hex{12340}.
\end{alertblock}
\Sumber{Bab VI, Bagian 1}
\end{frame}

% ---------------------------------------------------------
\begin{frame}{Rumus Alamat Fisik}
\small
\[
\text{alamat fisik} = (\text{segmen} \times 16) + \text{offset}
\]

\vspace{0.2em}
\begin{itemize}
    \item Pengali 16 sama dengan \hex{10}: segmen digeser empat bit
          ke kiri.
    \item Segmen menyatakan nomor paragraf, bukan nomor byte.
    \item Hasil dapat mencapai 20 bit karena segmen dan offset
          masing-masing 16 bit.
\end{itemize}

\vspace{0.2em}
Contoh:
\[
\text{1234h} \times 16 + \text{5678h}
= \text{12340h} + \text{5678h} = \text{179B8h}
\]
\Sumber{Bab VI, Bagian 1}
\end{frame}

% ---------------------------------------------------------
\begin{frame}{Pembentukan Alamat Fisik}
\begin{center}
\begin{tikzpicture}[
    digit/.style={
        draw=obeBlue, line width=0.5pt, minimum size=0.62cm,
        inner sep=0pt, font=\ttfamily\small, anchor=center
    },
    lab/.style={font=\scriptsize\bfseries, text=obeBlue, anchor=east}
]
    \foreach \d [count=\i from 0] in {1,2,3,4,0}
        \node[digit, fill=obeBlueLight] at (1.4+0.65*\i,2.2) {\d};
    \foreach \d [count=\i from 0] in {0,5,6,7,8}
        \node[digit, fill=obeGrayLight, draw=obeGray]
            at (1.4+0.65*\i,1.5) {\d};
    \foreach \d [count=\i from 0] in {1,7,9,B,8}
        \node[digit, fill=obeAmberLight, draw=obeAmber]
            at (1.4+0.65*\i,0.8) {\d};

    \node[lab] at (1.0,2.2) {Segmen};
    \node[lab] at (1.0,1.5) {Offset};
    \node[lab] at (1.0,0.8) {Fisik};
    \draw[line width=0.9pt, obeBlue] (1.05,1.12) -- (4.55,1.12);
\end{tikzpicture}
\end{center}

\vspace{0.2em}
{\footnotesize\color{obeGray}Digit 0 di kanan menandai pergeseran
empat bit; hasilnya \hex{179B8}.\par}
\Sumber{Bab VI, Bagian 1}
\end{frame}

% ---------------------------------------------------------
\begin{frame}{Enam Contoh Perhitungan}
\footnotesize
\begin{center}
\begin{tabular}{@{}>{\raggedright\arraybackslash}p{0.8cm}
    >{\raggedright\arraybackslash}p{1.6cm}
    >{\raggedright\arraybackslash}p{1.6cm}
    >{\raggedright\arraybackslash}p{3.9cm}
    >{\raggedright\arraybackslash}p{2.1cm}@{}}
    \toprule
    \textbf{No} & \textbf{Segmen} & \textbf{Offset} &
    \textbf{Perhitungan} & \textbf{Fisik} \\
    \midrule
    1 & \hex{0000} & \hex{0100} & \hex{00000} + \hex{0100} &
    \hex{00100} \\
    \addlinespace
    2 & \hex{1000} & \hex{0020} & \hex{10000} + \hex{0020} &
    \hex{10020} \\
    \addlinespace
    3 & \hex{076A} & \hex{0100} & \hex{076A0} + \hex{0100} &
    \hex{077A0} \\
    \addlinespace
    4 & \hex{2000} & \hex{FFFF} & \hex{20000} + \hex{FFFF} &
    \hex{2FFFF} \\
    \addlinespace
    5 & \hex{1234} & \hex{5678} & \hex{12340} + \hex{5678} &
    \hex{179B8} \\
    \addlinespace
    6 & \hex{B800} & \hex{0000} & \hex{B8000} + \hex{0000} &
    \hex{B8000} \\
    \bottomrule
\end{tabular}
\end{center}
\Sumber{Bab VI, Bagian 1}
\end{frame}

% ---------------------------------------------------------
\begin{frame}{Rangkuman Bagian 1}
\small
\begin{itemize}
    \item Register 16 bit hanya menjangkau 64 KB; bus alamat 20 bit
          menjangkau 1 MB.

    \vspace{0.15em}
    \item Alamat logis ditulis \eng{segment:offset}; alamat fisik
          adalah alamat 20 bit di bus alamat.

    \vspace{0.15em}
    \item Alamat fisik = (segmen $\times$ 16) + offset.

    \vspace{0.15em}
    \item Pengali 16 berasal dari paragraf 16 byte; segmen
          menyatakan nomor paragraf.

    \vspace{0.15em}
    \item Hasil perhitungan dapat memiliki lima digit
          heksadesimal, misalnya \hex{179B8}.
\end{itemize}
\end{frame}

% =========================================================
% BAGIAN 2
% =========================================================
\PemisahBagian{Paragraf 16 Byte, Segmen 64 KB, dan Overlapping}

% ---------------------------------------------------------
\begin{frame}{Paragraf 16 Byte}
\small
\begin{itemize}
    \item Memori dibagi menjadi satuan 16 byte yang disebut
          \term{paragraf}.

    \vspace{0.18em}
    \item Paragraf pertama: \hex{00000}--\hex{0000F}; paragraf
          kedua: \hex{00010}--\hex{0001F}.

    \vspace{0.18em}
    \item Nilai segmen adalah nomor paragraf, sehingga alamat dasar
          segmen selalu kelipatan 16 dan berakhir dengan digit 0.

    \vspace{0.18em}
    \item 65.536 paragraf $\times$ 16 byte = 1.048.576 byte = 1 MB.
\end{itemize}

\vspace{0.25em}
\begin{block}{Akibat}
\small Alamat dasar segmen tidak pernah berakhir dengan digit
selain 0.
\end{block}
\Sumber{Bab VI, Bagian 2}
\end{frame}

% ---------------------------------------------------------
\begin{frame}{Segmen 64 KB}
\small
\begin{itemize}
    \item Segmen adalah jendela 64 KB yang dimulai pada salah satu
          alamat paragraf.

    \vspace{0.18em}
    \item Offset 16 bit bergerak dari \hex{0000} sampai \hex{FFFF}.

    \vspace{0.18em}
    \item Bila dasar segmen adalah $B$, segmen mencakup alamat $B$
          sampai $B + \hex{FFFF}$, yaitu 65.536 byte.

    \vspace{0.18em}
    \item Jarak antara dua dasar segmen berurutan hanya 16 byte,
          sehingga jendela segmen saling menumpuk.
\end{itemize}

\vspace{0.25em}
\begin{block}{Overlapping}
\small Segmen \hex{0000} dan \hex{0001} hampir seluruhnya
bertumpuk; inilah \eng{overlapping} segmen.
\end{block}
\end{frame}

% ---------------------------------------------------------
\begin{frame}{Satu Alamat, Banyak Pasangan}
\footnotesize
\begin{center}
\begin{tabular}{@{}>{\raggedright\arraybackslash}p{0.8cm}
    >{\raggedright\arraybackslash}p{1.5cm}
    >{\raggedright\arraybackslash}p{1.5cm}
    >{\raggedright\arraybackslash}p{3.7cm}
    >{\raggedright\arraybackslash}p{2.6cm}@{}}
    \toprule
    \textbf{No} & \textbf{Segmen} & \textbf{Offset} &
    \textbf{Perhitungan} & \textbf{Fisik} \\
    \midrule
    1 & \hex{0000} & \hex{0050} & \hex{00000} + \hex{0050} &
    \hex{00050} \\
    \addlinespace
    2 & \hex{0001} & \hex{0040} & \hex{00010} + \hex{0040} &
    \hex{00050} \\
    \addlinespace
    3 & \hex{0002} & \hex{0030} & \hex{00020} + \hex{0030} &
    \hex{00050} \\
    \addlinespace
    4 & \hex{0003} & \hex{0020} & \hex{00030} + \hex{0020} &
    \hex{00050} \\
    \addlinespace
    5 & \hex{0004} & \hex{0010} & \hex{00040} + \hex{0010} &
    \hex{00050} \\
    \addlinespace
    6 & \hex{0005} & \hex{0000} & \hex{00050} + \hex{0000} &
    \hex{00050} \\
    \bottomrule
\end{tabular}
\end{center}
\Sumber{Bab VI, Bagian 2}
\end{frame}

% ---------------------------------------------------------
\begin{frame}{Mengapa Overlapping Terjadi}
\begin{center}
\begin{tikzpicture}[
    seg/.style={
        draw=obeBlue, line width=0.8pt, fill=obeBlueLight,
        rounded corners=2pt, minimum width=3.4cm,
        minimum height=0.9cm, align=center, font=\small
    },
    off/.style={
        draw=obeTeal, line width=0.8pt, fill=obeTealLight,
        rounded corners=2pt, minimum width=2.0cm,
        minimum height=0.75cm, align=center, font=\small\ttfamily
    },
    phys/.style={
        draw=obeAmber, line width=0.9pt, fill=obeAmberLight,
        rounded corners=2pt, minimum width=2.8cm,
        minimum height=1.2cm, align=center, font=\small
    },
    panah/.style={-{Stealth[length=2.2mm]}, line width=0.8pt, obeBlue}
]
    \node[seg] (s1) at (0,2.0)
        {Segmen \hex{0000}\\{\scriptsize dasar \hex{00000}}};
    \node[off] (o1) at (3.8,2.0) {\hex{0050}};
    \node[seg] (s2) at (0,0.6)
        {Segmen \hex{0001}\\{\scriptsize dasar \hex{00010}}};
    \node[off] (o2) at (3.8,0.6) {\hex{0040}};
    \node[phys] (f) at (7.6,1.3)
        {Alamat fisik\\\hex{00050}};

    \draw[panah] (s1) -- (o1);
    \draw[panah] (s2) -- (o2);
    \draw[panah] (o1.east) -- ++(0.5,0) -| (f.north);
    \draw[panah] (o2.east) -- ++(0.5,0) -| (f.south);
\end{tikzpicture}
\end{center}

\vspace{0.2em}
{\footnotesize\color{obeGray}Dua pasangan berbeda menunjuk lokasi
fisik yang sama.\par}
\Sumber{Bab VI, Bagian 2}
\end{frame}

% ---------------------------------------------------------
\begin{frame}{Berapa Banyak Pasangan?}
\small
\begin{itemize}
    \item Untuk alamat fisik $F$ di bawah \hex{10000}, banyaknya
          pasangan adalah $\lfloor F/16 \rfloor + 1$.

    \vspace{0.18em}
    \item Contoh \hex{00050}: $\lfloor \hex{50}/\hex{10}
          \rfloor + 1 = 5 + 1 = 6$ pasangan.

    \vspace{0.18em}
    \item Offset hanya bergerak 64 KB; setiap kenaikan segmen
          menggeser dasar 16 byte.

    \vspace{0.18em}
    \item Jumlah pasangan tidak pernah melebihi
          $65.536/16 = 4.096$.
\end{itemize}

\vspace{0.25em}
\begin{block}{Rata-rata}
\small $65.536 \times 65.536$ alamat logis berpadanan dengan
1.048.576 alamat fisik, sehingga satu alamat fisik rata-rata
memiliki 4.096 nama.
\end{block}
\end{frame}

% ---------------------------------------------------------
\begin{frame}{Operator OFFSET dan SEG}
\small
\begin{columns}[t]
\begin{column}{0.47\textwidth}
    {\color{obeTeal}\bfseries \code{OFFSET}}\\[0.25em]
    \begin{itemize}
        \item Menghasilkan jarak label dari awal segmen.
        \item Contoh: \code{MOV AX, OFFSET var\_byte} memberi
              \hex{0000}.
    \end{itemize}
\end{column}
\begin{column}{0.49\textwidth}
    {\color{obeMaroon}\bfseries \code{seg}}\\[0.25em]
    \begin{itemize}
        \item Menghasilkan nomor segmen label.
        \item Contoh: \code{MOV AX, seg var\_nilai} memberi
              \hex{076A}.
    \end{itemize}
\end{column}
\end{columns}

\vspace{0.4em}
\begin{alertblock}{Jangan tertukar}
\small \code{seg} memberi nomor segmen \hex{076A}, bukan alamat
dasar \hex{076A0}.
\end{alertblock}
\Sumber{Bab VI, Bagian 2}
\end{frame}

% ---------------------------------------------------------
\begin{frame}{Rangkuman Bagian 2}
\small
\begin{itemize}
    \item Memori dibagi menjadi paragraf 16 byte; nilai segmen
          adalah nomor paragraf.

    \vspace{0.15em}
    \item Satu segmen menjangkau 64 KB, dari offset \hex{0000}
          sampai \hex{FFFF}.

    \vspace{0.15em}
    \item \eng{Overlapping} membuat satu alamat fisik dapat
          dinyatakan oleh banyak pasangan.

    \vspace{0.15em}
    \item Banyaknya pasangan untuk $F$ adalah
          $\lfloor F/16 \rfloor + 1$, maksimum 4.096.

    \vspace{0.15em}
    \item \code{OFFSET} memberi jarak dari awal segmen;
          \code{seg} memberi nomor segmen.
\end{itemize}
\end{frame}

% =========================================================
% BAGIAN 3
% =========================================================
\PemisahBagian{Pembagian Memori 1 MB dan Verifikasi Alamat}

% ---------------------------------------------------------
\begin{frame}{Enam Belas Blok 64 KB}
\footnotesize
\begin{center}
\begin{tabular}{@{}>{\raggedright\arraybackslash}p{3.2cm}
    >{\raggedright\arraybackslash}p{1.8cm}
    >{\raggedright\arraybackslash}p{7.6cm}@{}}
    \toprule
    \textbf{Rentang alamat} & \textbf{Ukuran} &
    \textbf{Isi dan kegunaan} \\
    \midrule
    \hex{00000}--\hex{9FFFF} & 640 KB &
    Memori konvensional: program DOS, data, tabel vektor interupsi,
    data BIOS \\
    \addlinespace
    \hex{A0000}--\hex{BFFFF} & 128 KB &
    Buffer video untuk tampilan grafis dan teks \\
    \addlinespace
    \hex{C0000}--\hex{EFFFF} & 192 KB &
    ROM adaptor, termasuk BIOS kartu tambahan \\
    \addlinespace
    \hex{F0000}--\hex{FFFFF} & 64 KB &
    ROM BIOS sistem \\
    \bottomrule
\end{tabular}
\end{center}
\Sumber{Bab VI, Bagian 3}
\end{frame}

% ---------------------------------------------------------
\begin{frame}[fragile]{Verifikasi dengan DEBUG}
\begin{columns}[t]
\begin{column}{0.55\textwidth}
\begin{keluaran}
C:\>DEBUG hitung-alamat.exe
-R
DS=076A  CS=076A  IP=0100
-D DS:0000 L 08
076A:0000  CD 20 FF 9F 00 9A F0 FE
-Q
\end{keluaran}
\end{column}
\begin{column}{0.41\textwidth}
\footnotesize
Perintah \code{R} menampilkan register; \code{D} menampilkan isi
memori.

\vspace{0.3em}
\code{076A:0000} adalah alamat logis, bukan alamat fisik.

\vspace{0.3em}
Alamat fisiknya:
\[
(\hex{076A} \times 16) + \hex{0000} = \hex{076A0}
\]
\end{column}
\end{columns}

\vspace{0.2em}
{\footnotesize\color{obeGray}DEBUG menampilkan pasangan
\code{segmen:offset}; alamat fisik dihitung dengan rumus.\par}
\Sumber{Bab VI, Bagian 3}
\end{frame}

% ---------------------------------------------------------
\begin{frame}[fragile]{Membuktikan Overlapping dengan DEBUG}
\small
Dua pasangan berikut menunjuk alamat fisik \hex{076A0}:

\[
(\hex{0769} \times 16) + \hex{0010}
= \hex{07690} + \hex{0010} = \hex{076A0}
\]

\begin{columns}[t]
\begin{column}{0.55\textwidth}
\begin{keluaran}
-D 076A:0000 L 08
-D 0769:0010 L 08
\end{keluaran}
\end{column}
\begin{column}{0.41\textwidth}
\footnotesize
Kedua perintah menampilkan isi memori yang sama.
\end{column}
\end{columns}

\vspace{0.3em}
\begin{block}{Bukti}
\small Alamat logis berbeda, tetapi alamat fisiknya sama.
\end{block}
\Sumber{Bab VI, Bagian 3}
\end{frame}

% ---------------------------------------------------------
\begin{frame}{Alamat Logis vs Alamat Fisik di DEBUG}
\small
\begin{itemize}
    \item DEBUG menampilkan alamat dengan format
          \code{ssss:oooo}, yaitu segmen dan offset.

    \vspace{0.18em}
    \item Alamat fisik = \code{ssss} $\times 16$ + \code{oooo}.

    \vspace{0.18em}
    \item \code{076A:0000} dan \code{0769:0010} sama-sama menunjuk
          \hex{076A0}.

    \vspace{0.18em}
    \item \code{B800:0000} menunjuk \hex{B8000}, yaitu area buffer
          video.
\end{itemize}

\vspace{0.25em}
\begin{block}{Kebiasaan baik}
\small Saat menelusuri variabel dengan \eng{debugger}, catat
segmen dan offset sekaligus.
\end{block}
\end{frame}

% ---------------------------------------------------------
\begin{frame}{Rangkuman Bagian 3}
\small
\begin{itemize}
    \item Memori 1 MB dipandang sebagai enam belas blok 64 KB.

    \vspace{0.15em}
    \item Pembagiannya: memori konvensional 640 KB, buffer video
          128 KB, ROM adaptor 192 KB, dan ROM BIOS 64 KB.

    \vspace{0.15em}
    \item DEBUG menampilkan alamat logis \code{segmen:offset},
          bukan alamat fisik.

    \vspace{0.15em}
    \item Alamat fisik dihitung dengan rumus
          (segmen $\times$ 16) + offset.

    \vspace{0.15em}
    \item Dua pasangan berbeda dapat menampilkan isi memori yang
          sama; itulah bukti \eng{overlapping}.
\end{itemize}
\end{frame}

% =========================================================
% CONTOH PROGRAM
% =========================================================
\begin{frame}[fragile]{Contoh 1: Menampilkan Offset Variabel}
\begin{columns}[t]
\begin{column}{0.50\textwidth}
\begin{kodeasm}
.MODEL SMALL
.DATA
    var_byte DB 10h
    var_word DW 1234h
.CODE
START:
    MOV AX, @DATA
    MOV DS, AX
    MOV AX, OFFSET var_byte
    CALL cetak_hex
    MOV AX, OFFSET var_word
    CALL cetak_hex
...
\end{kodeasm}
\end{column}
\begin{column}{0.48\textwidth}
\footnotesize
{\color{obeTeal}\bfseries Keluaran}\\[0.15em]
\begin{keluaran}
Offset var_byte = 0000
Offset var_word = 0001
\end{keluaran}

\vspace{0.3em}
Bila \reg{DS} = \hex{076A}:
\begin{itemize}
    \item \code{var\_byte}: \hex{076A0} + \hex{0000} =
          \hex{076A0}
    \item \code{var\_word}: \hex{076A0} + \hex{0001} =
          \hex{076A1}
\end{itemize}

\vspace{0.2em}
{\color{obeGray}\code{OFFSET} memberi offset, bukan alamat fisik.
\par}
\end{column}
\end{columns}
\Sumber{Bab VI, Contoh 1}
\end{frame}

% ---------------------------------------------------------
\begin{frame}[fragile]{Contoh 2: Menghitung Alamat Fisik}
\begin{columns}[t]
\begin{column}{0.55\textwidth}
\begin{kodeasm}
.MODEL SMALL
.DATA
    var_nilai DW 0ABCDh
.CODE
START:
    MOV AX, @DATA
    MOV DS, AX
    MOV AX, seg var_nilai
    MOV BX, OFFSET var_nilai
    MOV DX, AX
    MOV CL, 4
    SHR DX, CL
    SHL AX, CL
    ADD AX, BX
    ADC DX, 0
\end{kodeasm}
\end{column}
\begin{column}{0.43\textwidth}
\footnotesize
{\color{obeTeal}\bfseries Keluaran}\\[0.15em]
\begin{keluaran}
Segmen: 076A
Offset: 0000
Alamat fisik: 000076A0
\end{keluaran}

\vspace{0.3em}
Hasil 20 bit disimpan pada \reg{DX}:\reg{AX}: DX memuat bit
16--19, AX memuat bit 0--15.

\vspace{0.2em}
{\color{obeGray}\code{SHL} = kali 16; \code{SHR} memindahkan bit
atas ke DX.\par}
\end{column}
\end{columns}
\Sumber{Bab VI, Contoh 2}
\end{frame}

% =========================================================
% PRAKTIKUM
% =========================================================
\begin{frame}{Praktikum: Menghitung dan Memverifikasi Alamat}
\begin{block}{Judul praktikum}
\small Menghitung dan memverifikasi alamat fisik dengan
\eng{debugger}\\[0.15em]
\textbf{Sub-CPMK:} 1.4 \quad\textbullet\quad
\textbf{Waktu:} 100 menit \quad\textbullet\quad
\textbf{Bentuk:} individu
\end{block}

\vspace{0.2em}
\begin{columns}[t]
\begin{column}{0.46\textwidth}
{\color{obeBlue}\bfseries Alat dan bahan}
\begin{itemize}\footnotesize
    \item Komputer dengan emu8086
    \item DOSBox dengan TASM, TLINK, dan DEBUG
    \item \file{alamat-fisik.asm} dan \file{hitung-alamat.asm}
    \item Kalkulator heksadesimal
\end{itemize}
\end{column}
\begin{column}{0.50\textwidth}
{\color{obeBlue}\bfseries Langkah kerja}
\begin{enumerate}\footnotesize
    \item Jalankan \file{alamat-fisik.asm}; catat offset.
    \item DEBUG \file{r}; catat DS; hitung alamat fisik.
    \item \file{D DS:0000 L 08}; bandingkan isi.
    \item Buktikan overlap dengan \code{0769:0011}.
    \item Jalankan \file{hitung-alamat.asm}; bandingkan hasil.
\end{enumerate}
\end{column}
\end{columns}

\vspace{0.2em}
{\footnotesize\textbf{Dikumpulkan:} perhitungan minimal enam
pasangan, tangkapan layar program dan DEBUG, bukti
\eng{overlapping}, serta jawaban Kuis 1.\par}
\Sumber{Bab VI, Praktikum}
\end{frame}

% =========================================================
% LATIHAN
% =========================================================
\begin{frame}{Latihan}
\small
\begin{enumerate}
    \item Jelaskan mengapa 8086 berbus alamat 20 bit memakai
          register segmen 16 bit. Apa yang hilang bila hanya satu
          register 16 bit dipakai?

    \vspace{0.15em}
    \item Hitung alamat fisik dan tuliskan langkahnya:
          \code{0000h:1A00h}, \code{1500h:0100h},
          \code{9FFFh:000Fh}, \code{B800h:0060h},
          \code{ABCDh:1234h}.

    \vspace{0.15em}
    \item Alamat fisik \hex{00080} dapat dinyatakan oleh beberapa
          pasangan. Tentukan banyaknya dan tuliskan seluruh
          pasangan.

    \vspace{0.15em}
    \item Sebutkan pembagian memori 1 MB dan kegunaan setiap
          wilayah.

    \vspace{0.15em}
    \item Seorang mahasiswa memperoleh \hex{68AC} untuk
          \code{1234h:5678h} dan menyimpulkan \hex{0000} dari
          tampilan \code{076A:0000}. Tunjukkan kesalahannya.
\end{enumerate}
\end{frame}

% =========================================================
% RANGKUMAN
% =========================================================
\begin{frame}{Rangkuman}
\small
\begin{itemize}
    \item Segmentasi menjembatani register 16 bit dan bus alamat
          20 bit.

    \vspace{0.15em}
    \item Alamat logis adalah pasangan \eng{segment:offset};
          alamat fisik adalah alamat 20 bit di bus alamat.

    \vspace{0.15em}
    \item Alamat fisik = (segmen $\times$ 16) + offset.

    \vspace{0.15em}
    \item Paragraf berukuran 16 byte; satu segmen menjangkau
          64 KB; enam belas segmen menutupi 1 MB.

    \vspace{0.15em}
    \item \eng{Overlapping} membuat satu alamat fisik memiliki
          banyak pasangan, maksimum 4.096.

    \vspace{0.15em}
    \item DEBUG menampilkan alamat logis; alamat fisik dihitung
          dengan rumus.
\end{itemize}
\end{frame}

% ---------------------------------------------------------
\begin{frame}{Cek Pemahaman}
\footnotesize
\begin{enumerate}
    \item Rumus yang benar untuk alamat fisik adalah \dots
          \begin{enumerate}
              \item segmen + offset
              \item (segmen $\times$ \hex{10}) + offset
              \item (segmen $\times$ 20) + offset
              \item (offset $\times$ \hex{10}) + segmen
          \end{enumerate}

    \vspace{0.2em}
    \item Nilai segmen \hex{1234} menunjuk alamat dasar \dots
          \begin{enumerate}
              \item \hex{1234} \quad \item \hex{12340} \quad
              \item \hex{01234} \quad \item \hex{23400}
          \end{enumerate}

    \vspace{0.2em}
    \item Alamat fisik yang ditunjuk \code{B800h:0000h} adalah
          \dots
          \begin{enumerate}
              \item \hex{B800} \quad \item \hex{0B800} \quad
              \item \hex{B8000} \quad \item \hex{BB000}
          \end{enumerate}
\end{enumerate}
\end{frame}

% =========================================================
% PENUTUP
% =========================================================
\begin{frame}{Penutup}
\begin{block}{Sebelum pertemuan berikutnya}
\small
\begin{itemize}
    \item Selesaikan praktikum: perhitungan alamat, tangkapan layar
          DEBUG, bukti \eng{overlapping}, dan Kuis 1.
    \item Baca \textbf{Bab VII} — Memulai Assembly: TASM, TLINK,
          COM, EXE, dan Direktif.
    \item Siapkan pertanyaan tentang struktur program dan alur
          translasi.
\end{itemize}
\end{block}

\vspace{0.4em}
\begin{center}
{\color{obeBlue}\sffamily\large\bfseries Terima kasih\par}
\vspace{0.15em}
{\color{obeGray}\sffamily\small \NamaMK\ (\KodeMK)
\textbullet\ \SemesterTeks\par}
\end{center}
\end{frame}

\end{document}
